% Lâminas de motivação: SC26 (Supercomputing 2026) — Relevância e Escala.
% Aula 11 (NumPy): transforma a aula prática em uma aula de 1 hora.
\documentclass[aspectratio=169, 10pt]{beamer}
\usepackage{beamer_theme_cienciadados2026}

\title[SC26 — Supercomputing 2026]{Supercomputing 2026 (SC26): Relevância e Escala}
\subtitle{Motivação para a Aula 11 — Computação Científica com NumPy}
\author{Prof. Dr. Ney Lemke}
\institute{UNESP -- Instituto de Biociências de Botucatu}
\date{Edição 2026 · a partir da participação no SC26}

\begin{document}

\frame{\titlepage}

\begin{frame}{Por que começar por aqui?}
  \begin{block}{Da sua máquina ao topo do mundo da computação}
    Nesta aula vamos do \textbf{NumPy} (nosso primeiro degrau) até a \textbf{supercomputação}:
    a escala de máquinas e problemas que movem a ciência moderna. O SC26 é o palco onde
    essa comunidade se encontra.
  \end{block}
  \begin{itemize}
    \item Entender \textbf{onde} a computação científica chegou (exascale).
    \item Ver a \textbf{relevância} do HPC para a ciência e para a IA.
    \item Conectar isso à ferramenta que vamos praticar: o \textbf{NumPy}.
  \end{itemize}
\end{frame}

\begin{frame}{Sumário das Lâminas}
  \tableofcontents
\end{frame}

\section{O Evento}

\begin{frame}{O que é o SC?}
  \begin{block}{A maior conferência mundial de supercomputação}
    \textbf{SC — The International Conference for High Performance Computing, Networking,
    Storage, and Analysis}, o evento anual mais importante da comunidade de HPC
    (patrocinado por ACM/SIGHPC e IEEE Computer Society).
  \end{block}
  \begin{itemize}
    \item \textbf{38ª edição}, série que acontece desde 1988.
    \item \textbf{SC26}: 15--20 de novembro de 2026, \textbf{McCormick Place, Chicago}.
    \item Tema da edição: \textbf{``HPC Unites''} (o HPC une comunidades e ciências).
    \item Exposição, programa técnico revisado por pares, \emph{Students@SC},
          e o famoso \textbf{SCinet} (a maior rede temporária do mundo, montada para o evento).
  \end{itemize}
\end{frame}

\begin{frame}{A escala de um evento gigante}
  \begin{columns}
    \begin{column}{0.5\textwidth}
      \begin{block}{Participantes}
        $\sim$ \textbf{18 a 20 mil} pesquisadores, engenheiros e estudantes de todo o mundo.
      \end{block}
      \begin{block}{Exposição}
        \textbf{500+ empresas} e instituições \\ (laboratórios nacionais, universidades, startups).
      \end{block}
    \end{column}
    \begin{column}{0.48\textwidth}
      \begin{block}{Dias / área}
        6 dias de programação \\ (exposição em 3 dias) no \textbf{McCormick Place}, à beira do Lago Michigan.
      \end{block}
      \begin{alertblock}{IEEE 80 anos}
        A IEEE Computer Society celebrou \textbf{80 anos} nesta edição.
      \end{alertblock}
    \end{column}
  \end{columns}
\end{frame}

\section{Escala de Computação}

\begin{frame}{A escala numérica da computação}
  \begin{block}{Operações de ponto flutuante por segundo (FLOPS)}
    \centering
    \begin{tabular}{l l l}
      \toprule
      \textbf{Sistema} & \textbf{Desempenho} & \textbf{Ordem} \\
      \midrule
      Seu laptop & $\sim 10^{11}$--$10^{13}$ & giga--tera \\
      Cluster de laboratório & $\sim 10^{15}$ & peta \\
      \textbf{Máquinas exascale} & $\sim 10^{18}$ & \textbf{exa} \\
      Zettascale (próxima fronteira) & $\sim 10^{21}$ & zetta \\
      \bottomrule
    \end{tabular}
  \end{block}
  \begin{itemize}
    \item O exascale é \textbf{um milhão de vezes} mais rápido que um laptop comum.
    \item Frontier (ORNL), Aurora (ANL) e El Capitan (LLNL) são as máquinas de ponta nos EUA.
  \end{itemize}
\end{frame}

\begin{frame}{Por que tamanha rapidez importa?}
  \begin{columns}
    \begin{column}{0.5\textwidth}
      \begin{block}{Ciência}
        \begin{itemize}
          \item Clima e tempo.
          \item Descoberta de materiais e fármacos.
          \item Física fundamental e fusão.
          \item Simulação de proteínas e genômica.
        \end{itemize}
      \end{block}
    \end{column}
    \begin{column}{0.48\textwidth}
      \begin{block}{IA em escala}
        \begin{itemize}
          \item Treinar modelos gigantes exige \textbf{exascale}.
          \item HPC e IA convergiram; viraram um só problema.
          \item ``Agentic AI'' para ciência: agentes que planejam e rodam experimentos.
        \end{itemize}
      \end{block}
    \end{column}
  \end{columns}
\end{frame}

\section{O que Marcou o SC26}

\begin{frame}{Temas-chave da edição 2026}
  \begin{alertblock}{Keynote (Anima Anandkumar, Caltech)}
    ``\textbf{The Next Scaling Law after LLMs: AI that Understands the Physical World}''
    --- a ideia: a próxima grande evolução vem de IA treinada em \textbf{dados do mundo físico},
    não apenas em texto.
  \end{alertblock}
  \begin{itemize}
    \item \textbf{HPC Unites}: convergência entre supercomputação, IA e ciência.
    \item \textbf{Exascale $\to$ Zettascale}: os próximos mil passos de capacidade.
    \item \textbf{Agentes como cientistas} (NVIDIA: \emph{agents as scientists}), laboratórios autônomos.
    \item Computação \textbf{quântica} entrando no datacenter de HPC.
  \end{itemize}
\end{frame}

\begin{frame}{Python e NumPy no coração do HPC}
  \begin{block}{Surpresa boa da conferência}
    No SC26 houve o workshop \textbf{``High Performance Python for Science at Scale''}
    e falas sobre \textbf{Python como ``tecido conjuntivo''} de workflows exascale.
  \end{block}
  \begin{itemize}
    \item Até na supercomputação, \textbf{Python/NumPy} é a camada que une simulação, IA e análise.
    \item O que você aprende hoje (\textbf{arrays vetorizados}) é o \textbf{mesmo princípio} que roda nas máquinas exascale.
    \item Vetorização, broadcasting e reduções \textbf{não são só ``academia''}: são otimização a sério.
  \end{itemize}
\end{frame}

\begin{frame}{Momento SC26}
  \begin{block}{Registros do evento}
    Espaço para fotos da sua participação (stand dos laboratórios, exposição, Chicago).
  \end{block}
  \begin{enumerate}
    \item Impressão da \textbf{escala} (número de participantes e expositores).
    \item \textbf{Demonstrações} de IA científica e simulação em tempo real.
    \item \textbf{Conexão}: ver ``de perto'' que computação de alta performance move a ciência.
  \end{enumerate}
\end{frame}

\begin{frame}{Resumo da Aula}
  \begin{itemize}
    \item \textbf{SC26}: o maior evento de HPC do mundo (Chicago, 38ª edição, ``HPC Unites'').
    \item \textbf{Escala}: exascale ($10^{18}$ FLOPS) é $\sim 10^{6}$ vezes um laptop.
    \item \textbf{Relevância}: ciência e IA dependem de HPC; as fronteiras (zettascale, IA do mundo físico, agentes) estão abertas.
    \item \textbf{NumPy} é a base do ecossistema e a nossa prática na Aula 11.
    \item \textbf{Prática}: criar, conhecer, indexar, filtrar, vetorizar e resumir arrays.
  \end{itemize}
\end{frame}

\end{document}